\documentclass[11pt]{article}
\usepackage[margin=1in]{geometry}
\usepackage{amsmath,amssymb,booktabs}

\title{Depth Shortcut Predictor with NormEncoder Input: Training Summary}
\date{}

\begin{document}
\maketitle

\section{Goal}
This branch keeps the predictor in direction-normalized space while giving it
explicit source-layer magnitude context. The shortcut predictor now receives
both the token direction \(U_a\) and a small encoded log-norm map \(E_m\):
\[
    P_\theta(U_a,m_a,a,b,t)
    \rightarrow
    (Y_{a\rightarrow b}, \hat{\delta m}_{a\rightarrow b}^{\mathrm{norm}}).
\]

\section{Source State Inputs}
For a source hidden state
\[
    Z_a\in\mathbb{R}^{B\times256\times768},
\]
compute tokenwise magnitude, direction, and log magnitude:
\[
    r_a(i)=\|Z_a(i)\|_2,
    \qquad
    U_a(i)=\frac{Z_a(i)}{r_a(i)+10^{-6}},
\]
\[
    m_a(i)=\ln(r_a(i)+10^{-6}).
\]
The direction \(U_a\) remains the main representation input.

\section{Norm Map}
For SiT-B/2, build two log-norm channels:
\[
    M_a^{\mathrm{abs}}(i)=\frac{m_a(i)-5.5}{1.5},
\]
\[
    M_a^{\mathrm{rel}}(i)=
    \frac{m_a(i)-\operatorname{mean}_{j}m_a(j)}
    {\operatorname{std}_{j}m_a(j)+10^{-6}}.
\]
Concatenate and reshape:
\[
    M_a=[M_a^{\mathrm{abs}},M_a^{\mathrm{rel}}]
    \in\mathbb{R}^{B\times16\times16\times2}.
\]
The map is stop-gradient before entering the predictor:
\[
    M_a^{\mathrm{in}}=\operatorname{sg}(M_a).
\]

\section{NormEncoder}
The norm map is encoded by a small local encoder:
\[
    E_m=\mathrm{NormEncoder}(M_a^{\mathrm{in}}),
    \qquad
    E_m\in\mathbb{R}^{B\times16\times16\times32}.
\]
The current implementation is
\[
    \mathrm{Conv}_{1\times1}(2\to32)
    \rightarrow
    \mathrm{GELU}
    \rightarrow
    \mathrm{DWConv}_{3\times3}(32)
    \rightarrow
    \mathrm{GELU}.
\]

\section{Predictor Input}
Reshape the direction into a grid:
\[
    U_a^{\mathrm{grid}}\in\mathbb{R}^{B\times16\times16\times768}.
\]
Concatenate the direction grid and norm encoding:
\[
    X_a=[U_a^{\mathrm{grid}},E_m]
    \in\mathbb{R}^{B\times16\times16\times800}.
\]
Then project to predictor width \(C\):
\[
    h_0=\mathrm{Linear}(800\to C)(X_a).
\]
Examples:
\[
    C=256\ \text{for Tiny},\qquad
    C=384\ \text{for Small},\qquad
    C=768\ \text{for Base/Large}.
\]

\section{Conditioning}
The depth and timestep condition is unchanged:
\[
    h_{\mathrm{depth}}=[h_a,h_b,h_b-h_a,h_\Delta],
\]
\[
    c=\mathrm{MLP}\left(
    [h_{\mathrm{depth}},\operatorname{sg}(e_t^{\mathrm{DiT}})]
    \right).
\]
This condition modulates the predictor backbone through AdaLN/FiLM.

\section{Backbone and Heads}
The projected input \(h_0\) is processed by the configured predictor backbone:
\[
    h=\mathrm{Backbone}(h_0,c).
\]
The branch supports the existing ConvNeXt family and the architecture
ablations:
\[
    \texttt{convnext\_*},\qquad
    \texttt{dilated\_*},\qquad
    \texttt{attn\_hybrid\_*}.
\]

The direction head predicts a residual update:
\[
    Y_{a\rightarrow b}=U_a+\gamma_{\mathrm{out}}\Delta Y.
\]
The direction used for loss or skip inference is
\[
    \hat U_{a\rightarrow b}=\operatorname{dir}(Y_{a\rightarrow b}).
\]

In direction--magnitude modes, the magnitude head receives \([h,M_a]\) and
predicts a bounded normalized log-magnitude delta:
\[
    \hat{\delta m}_{a\rightarrow b}^{\mathrm{norm}}=\tanh(o_m).
\]
With SiT-B/2 scale \(s_m=3\),
\[
    \hat{\delta m}_{a\rightarrow b}
    =
    3\,\hat{\delta m}_{a\rightarrow b}^{\mathrm{norm}}.
\]

\section{Training Loss}
Each iteration still runs a full DiT forward to obtain
\[
    Z_0,\ldots,Z_L,\qquad \hat v.
\]
The total loss in direction--magnitude mode is
\[
    \mathcal{L}
    =
    \mathcal{L}_{\mathrm{FM}}
    +\lambda_{\mathrm{dir}}\mathcal{L}_{\mathrm{dir}}
    +\lambda_{\mathrm{boot}}\mathcal{L}_{\mathrm{boot}}
    +\lambda_{\mathrm{mag}}\mathcal{L}_{\mathrm{mag}}
    +\lambda_{\mathrm{boot-mag}}\mathcal{L}_{\mathrm{boot-mag}}.
\]
Default weights:
\[
    \lambda_{\mathrm{dir}}=0.5,\quad
    \lambda_{\mathrm{boot}}=0.25,\quad
    \lambda_{\mathrm{mag}}=0.375,\quad
    \lambda_{\mathrm{boot-mag}}=0.1875.
\]

\section{Skip-in-the-Loop Mode}
In \texttt{direction-magnitude-skip} mode, a stochastic skip-FM branch is added
after warmup:
\[
    z_{\mathrm{skip}}\sim\mathrm{Bernoulli}(p_{\mathrm{skip}}).
\]
When active, sample \(a<b\), restore
\[
    \tilde Z_b
    =
    \exp\left(
    \operatorname{clip}
    (m_a+3\hat{\delta m}_{a\rightarrow b}^{\mathrm{norm}},3,8)
    \right)
    \hat U_{a\rightarrow b},
\]
resume the DiT from layer \(b+1\), and compute
\[
    \mathcal{L}_{\mathrm{skip-FM}}
    =
    \|\hat v_{\mathrm{skip}}-(x_0-x_1)\|_2^2.
\]
The total loss becomes
\[
    \mathcal{L}
    \leftarrow
    \mathcal{L}
    +\lambda_{\mathrm{skip-FM}}\mathcal{L}_{\mathrm{skip-FM}}.
\]

\section{Direction-Only Compatibility}
For direction-only skip inference or skip-FID, the predictor still receives
the real source log magnitude \(m_a\), so NormEncoder is active. The caller
sets \texttt{return\_magnitude=False}, and only \(Y_{a\rightarrow b}\) is
returned. In direction-only mode, target magnitudes are still restored by the
stored L2 EMA calibration rather than by the magnitude head.

\end{document}
